\section{Discussion}\label{sec:discussion}
\paragraph{What is proved, what is designed, what is measured} \Cref{tab:status} separates the three kinds of claim. The
theorems hold under their hypotheses; the algorithm uses them as design rules, not as guarantees about its output; and
the numbers are evidence on synthetic data.

\paragraph{Limits that no amount of data removes} Some limits are theorems, and the output is built to show them rather
than hide them. \emph{Error floor}: in the Gaussian model of \cref{sec:group2} every unbiased estimator of the components
has variance at least $\noise^2/\gN$ in the worst direction, so closely aligned groups are intrinsically hard to separate.
\emph{No data-free threshold}: whether a numerical symbol ``factors'' depends on the noise level, which is why decisions are
made by evidence and not by a tolerance on coefficients. \emph{Repeated factors}: finite noisy data cannot distinguish
$P_i^2$ from two nearby factors, hence the label \UND. \emph{No certificate}: by \cref{prop:certify}, outputs are
compatibility statements with uncertainty, never identifications.

\paragraph{Beyond the linear ring} The two groupings do not depend on the equation being linear with constant
coefficients in one space dimension, and there are concrete reasons to expect the algorithm to extend. The matrix pencil
reads the temporal rates at every wavenumber vector $\kappa\in\R^d$ exactly as in one dimension, and the groups become
hypersurfaces $Q_j(\lambda,i\kappa)=0$ that grouping I can fit in the same way. For slowly varying coefficients a windowed
transform gives local rates $\lambda(\kappa;x_0)$, and each group becomes a family of curves indexed by position. For a
nonlinear equation, a short-time expansion shows that the flow is its linearization plus a quadratic interaction term,
so modes of the linearized flow can still be grouped on short windows. What does not carry over unchanged is the
theory: in higher dimension two factors can share a whole curve of zeros, the operators with variable coefficients form a
non-commutative algebra without unique factorization, and nonlinear solutions do not superpose. The first task in each
direction is therefore the one carried out here for the linear case: decide which pieces of the equation a solution
carries, how uniquely and how stably, and only then design the learning procedure around them (\SMbeyond).

\section{Conclusion}\label{sec:conclusion}
Groups are factors. This one statement, made precise by \cref{thm:factor} and its companions, determines what an
equation-learning algorithm should look for (modes on the curves of the factors), how it should combine evidence
(clustering factors across posterior samples), how it should measure uncertainty (through the Friedrichs angle), and
what it may claim (compatibility, never certification). On $285$ random equations never used in its development the
resulting method is right or equivalent on $96.5\%$ and never wrong without warning, while regression on a library of
terms is wrong without warning on most products and higher-order factors.

\section*{Code and data availability}
The code, the benchmark scripts and every stored result behind the numbers of this paper are at
\url{https://github.com/FlavioApric/pde-factor-fit}~\cite{Pri26code}; a snapshot is included in the supplementary
materials.

